\exercisesection{}

仅在习题 1 至 5 中，K 表示一个完备的非阿基米德、非离散的赋值域，其赋值记为 \(x \mapsto |x|\) 。我们还用 \(\mathrm{G} \subset \mathbf{R}_{+}^{*}\) 表示 \(\mathrm{K}^{*}\) 在映射 \(x \mapsto |x|\) 下的像。所考虑的巴拿赫空间都是 \(\mathrm{K}\) 上的巴拿赫空间。

\begin{enumerate}
\def\labelenumi{\arabic{enumi})}
\tightlist
\item
  \begin{enumerate}
  \def\labelenumii{\alph{enumii})}
  \tightlist
  \item
    设 E 是 K 上的一个巴拿赫空间。E 上存在一个范数 \(x \mapsto \|x\|\)，它给出 E 的拓扑，且使得 \(\|x\| \in \overline{G}\) 对所有 \(x \in E\) 成立。这样的范数称为圆范数。若 K 的赋值不是离散的，则每个给出 E 的拓扑的范数都是圆范数。
  \end{enumerate}
\end{enumerate}

\begin{enumerate}
\def\labelenumi{\alph{enumi})}
\setcounter{enumi}{1}
\item
  设 E 和 F 是 K 上的巴拿赫空间。我们在空间 \(\mathcal{L}(\mathrm{E};\mathrm{F})\) 上赋予范数 \(\| u\| = \sup_{x\in \mathrm{E}- \{0\}}\| u(x)\| /\| x\|\) 。空间 \(\mathcal{L}(\mathrm{E};\mathrm{F})\) 是 K 上的巴拿赫空间。若 E 上的范数是圆范数，则 \(\| u\| = \sup_{\| x\| \leqslant 1}\| u(x)\|\) 。若 F 上的范数也是圆范数，则 \(\mathcal{L}(\mathrm{E};\mathrm{F})\) 上的范数是圆范数。我们记 \(\mathrm{E}' = \mathcal{L}(\mathrm{E};\mathrm{K})\)，并称 \(\mathrm{E}'\) 为 E 的对偶空间。
\item
  设 I 是一个集合，F 是 K 上的一个巴拿赫空间，\(\mathcal{C}_{0}(\mathrm{I},\mathrm{F})\) 是由函数 \(f\colon I\to F\) 组成的 K-向量空间，其中 \(\|f(i)\|\) 沿 I 的有限子集的余集滤子趋于 0。赋以范数 \(\|f\|=\sup_{i\in I}\|f(i)\|\) 后，空间 \(\mathcal{C}_{0}(\mathrm{I},\mathrm{F})\) 是 K 上的巴拿赫空间，且若 F 的范数是圆范数，则它的范数也是圆范数。我们记 \(\mathcal{C}_{0}(\mathrm{I})=\mathcal{C}_{0}(\mathrm{I},\mathrm{K})\) 。
\end{enumerate}

\(d\) ) 设 I 是一个集合。对每个 \(i\in \mathrm{I}\) ，用 \(\varphi_{i}\in \mathcal{C}_{0}(\mathrm{I})\) 表示 \(\{i\}\) 的特征函数。设 F 是 K 上的一个巴拿赫空间。映射 \(u \mapsto (u(\varphi_i))_{i \in \mathrm{I}}\) 是从巴拿赫空间 \(\mathcal{L}(\mathcal{C}_0(\mathrm{I});\mathrm{F})\) 到巴拿赫空间 \(\mathcal{B}(\mathrm{I};\mathrm{K})\) 的同构（EVT, I, p. 4, 例；参见 EVT, I, p. 23, 习题 7）。

\begin{enumerate}
\def\labelenumi{\alph{enumi})}
\setcounter{enumi}{4}
\item
  假设 K 的赋值是离散的。设 E 是 K 上的范数为圆范数的巴拿赫空间。存在一个集合 I，使得 E 等距同构于巴拿赫空间 \(\mathcal{C}_{0}(\mathrm{I})\)（参见 EVT, I, p. 23, 习题 7）。
\item
  假设 K 的赋值是离散的。设 E 是赋有圆范数的 K 上的巴拿赫空间，F 是 E 的一个闭向量子空间。存在 E 的一个连续投影，其像为 F 且范数 \(\leqslant 1\) 。
\end{enumerate}

\begin{enumerate}
\def\labelenumi{\arabic{enumi})}
\setcounter{enumi}{1}
\tightlist
\item
  设 E 和 F 是 K 上的巴拿赫空间。我们称 \(u \in \mathcal{L}(\mathrm{E};\mathrm{F})\) 是完全连续的，如果它属于 \(\mathcal{L}(\mathrm{E};\mathrm{F})\) 中由有限秩线性映射组成的空间 \(\mathcal{L}^{\mathrm{f}}(\mathrm{E};\mathrm{F})\) 的闭包。我们用 \(\mathcal{L}^{\mathrm{c}}(\mathrm{E};\mathrm{F})\) 表示从 E 到 F 的完全连续线性映射组成的空间，并令 \(\mathcal{L}^{\mathrm{c}}(\mathrm{E}) = \mathcal{L}^{\mathrm{c}}(\mathrm{E};\mathrm{E})\) 。 a) 空间 \(\mathcal{L}^{\mathrm{c}}(\mathrm{E})\) 是 \(\mathcal{L}(\mathrm{E})\) 的一个闭双边理想。设 \(u \in \mathcal{L}(\mathrm{E};\mathrm{F})\) 。映射 \(\ell \mapsto \ell \circ u\) 是从 \(F'\) 到 \(E'\) 的连续线性映射，记为 \({}^{t}u\) 。若 \(u \in \mathcal{L}^{\mathrm{c}}(\mathrm{E};\mathrm{F})\) ，则 \({}^{t}u \in \mathcal{L}^{\mathrm{c}}(\mathrm{F}'; \mathrm{E}')\) 。
\end{enumerate}

\begin{enumerate}
\def\labelenumi{\alph{enumi})}
\setcounter{enumi}{1}
\item
  设 J 是一个集合，\(\mathrm{F} = \mathcal{C}_{0}(\mathrm{J})\) 。设 \(u \in \mathcal{L}(\mathrm{E};\mathrm{F})\) 。对每个 \(j \in J\)，映射 \(\ell_{j} \colon x \mapsto u(x)(j)\) 是 \(E'\) 的一个元素。我们有 \(\|u\| = \sup_{j \in J} \| \ell_{j} \|\)，并且映射 \(u \mapsto (\ell_{j})_{j \in J}\) 通过过渡到子空间诱导一个从 \(\mathcal{L}^{\mathrm{c}}(\mathrm{E};\mathrm{F})\) 到巴拿赫空间 \(\mathcal{C}_{0}(\mathrm{J};\mathrm{E}')\) 的同构。
\item
  若 K 是局部紧的，则 \(u \in \mathcal{L}(\mathrm{E}; \mathrm{F})\) 是完全连续的，当且仅当 E 的每个有界子集在 u 下的像在 F 中是相对紧的。（可将此结果与 III, p. 16 的命题 11 相比较。）
\end{enumerate}

\begin{enumerate}
\def\labelenumi{\arabic{enumi})}
\setcounter{enumi}{2}
\tightlist
\item
  我们用 A 表示 K 的由元素 \(x \in \mathrm{K}\)（满足 \(|x| \leqslant 1\)）组成的子环；它是一个局部环，其极大理想由 \(x \in \mathrm{A}\)（满足 \(|x| < 1\)）组成。
\end{enumerate}

设 I 是一个集合，E 是巴拿赫空间 \(\mathcal{C}_{0}(\mathrm{I})\)（参见习题 1 的问题 c)）。对 \(i \in I\)，用 \(\varphi_{i} \in E\) 表示 \(\{i\}\) 的特征函数。设 \(u \in \mathcal{L}^{\mathrm{c}}(\mathrm{E})\) 满足 \(\|u\| \leqslant 1\) 。

\begin{enumerate}
\def\labelenumi{\alph{enumi})}
\item
  子空间 \(E_{0}\) 由这样的元素 \(x \in E\) 组成：\(\|x\| \leqslant 1\)；它在 u 下稳定。
\item
  对 A 的每个非零理想 a，线性映射 u 通过过渡到商定义一个自同态 \(u_{a}\)，它作用于 A/a-模 \(E_{a} = E_{0}/aE_{0}\) 。这个 A/a-模是自由的，且 \(u_{a}\) 的像包含于一个有限生成的子模中。
\item
  存在唯一的形式级数 \(f_{u} \in A[[t]]\)，使得对 A 的每个非零理想 a，\(f_{u}\) 在 \((A/\mathfrak{a})[[t]]\) 中的像等于 A, VIII, p. 455 中定义的多项式 \(\det(1 - tu_{\mathfrak{a}})\)。
\end{enumerate}

\(d)\) 设 \(v \in \mathcal{L}^{\mathrm{c}}(\mathrm{E})\)，并设 \(c \in \mathrm{K}^{*}\) 满足 \(\|cv\| \leqslant 1\) 。形式级数 \(f_{cv}\) 与满足此性质的 \(c\) 的选取无关。

我们将用 \(\det(1-tv)\) 表示形式级数 \(f_v\)，这里 \(c \in \mathbf{K}^*\) 为任意满足 \(\|cv\| \leqslant 1\) 的元素。

¶ 4) 我们沿用上一习题的记号。

\begin{enumerate}
\def\labelenumi{\alph{enumi})}
\item
  设 \(v \in \mathcal{L}^c(\mathrm{E})\) 。形式级数 \(\det(1 - tv)\) 的收敛半径是无穷大。
\item
  设 \((v_{n})_{n\in\mathbf{N}}\) 是 \(\mathcal{L}^{c}(\mathrm{E})\) 中收敛于 \(v\in\mathcal{L}^{c}(\mathrm{E})\) 的一个序列。形式级数序列 \(\det(1-tv_{n})\) 对于形式级数系数的逐点收敛拓扑收敛于 \(\det(1-tv)\)。
\item
  设 \(v \in \mathcal{L}^{c}(E)\) 。对 K 的每个扩张 L 和每个 \(x \in L\)，我们用 \(\det(1 - xv)\) 表示形式幂级数 \(\det(1 - tv)\) 在 x 处的值，它由 a) 良定义。我们有 \(\det(1 + u)\det(1 + v) = \det(1 + u + v + u \circ v)\) 。
\item
  设 J 是一个集合，\(\mathrm{F} = \mathcal{C}_{0}(\mathrm{J})\) 。对所有 \(u \in \mathcal{L}^{\mathrm{c}}(\mathrm{E};\mathrm{F})\) 和 \(v \in \mathcal{L}^{\mathrm{c}}(\mathrm{F};\mathrm{E})\)，作为形式幂级数有 \(\det(1 - tu \circ v) = \det(1 - tv \circ u)\) 。
\item
  设 \(v \in \mathcal{L}^{\mathrm{c}}(\mathrm{E})\) 。我们用 \(\operatorname{Tr}(v)\) 表示 \(\det(1 - tu)\) 中 t 的系数的相反数。若 K 的特征为零，则
\end{enumerate}

\[
\det (1 - t u) = \exp \left(- \sum_ {m \geqslant 1} \frac {\operatorname{Tr} \left(u ^ {m}\right)}{m} t ^ {m}\right)
\]

作为形式幂级数成立。

¶ 5) 我们保留上一习题的记号。设 \(v \in \mathcal{L}^{\mathrm{c}}(\mathrm{E})\) 。

\begin{enumerate}
\def\labelenumi{\alph{enumi})}
\item
  形式幂级数 \(\det(1 - tv)\sum_{k\geqslant 0}t^{k}v^{k}\)（\(\mathcal{L}(\mathrm{E})[[t]]\) 中的）的收敛半径为无穷大。
\item
  设 \(\lambda \in \mathrm{K}\) 。E 的自同态 \(1_{\mathrm{E}} - \lambda v\) 可逆，当且仅当 \(\det(1 - \lambda v) \neq 0\) 。
\item
  设 \(\lambda \in K\) 满足 \(\det(1 - \lambda v) = 0\) 。存在 E 的唯一的闭子空间 \(I_{\lambda}\) 和 \(N_{\lambda}\)，它们在 \(v\) 下稳定，使得
\end{enumerate}

\begin{enumerate}
\def\labelenumi{(\roman{enumi})}
\item
  空间 E 是 \(I_{\lambda}\) 与 \(N_{\lambda}\) 的拓扑直和；
\item
  \(I_{\lambda}\) 的由 \(1_{E} - \lambda v\) 通过过渡到子空间得到的自同态是可逆的；
\item
  \(N_{\lambda}\) 的由 \(1_{E} - \lambda v\) 通过过渡到子空间得到的自同态是幂零的。
\end{enumerate}

\(d)\) 设 \(h \geqslant 1\) 是 det(1 - tv) 的零点 \(\lambda\) 的阶。我们有 \(\dim(N_{\lambda}) = h\) 。

应将本习题的结果与 III, p. 77 的 no. 5 中关于 R 或 C 上巴拿赫空间的结果相比较。

在下文中，K 表示域 R 或 C，且所有拓扑向量空间都是 K 上的。

\begin{enumerate}
\def\labelenumi{\arabic{enumi})}
\setcounter{enumi}{5}
\tightlist
\item
  设 E 是空间 \(C^{N}\)，赋以与半范数序列 \((p_{n})_{n\in\mathbf{N}}\) 相伴的拓扑，其中 \(p_{n}(x)=|x_{n}|\) 对 E 的每个元素 \(x=(x_{n})_{n\in\mathbf{N}}\) 成立。
\end{enumerate}

\begin{enumerate}
\def\labelenumi{\alph{enumi})}
\item
  空间 E 是一个弗雷歇空间。
\item
  E 的紧自同态组成的集合 \(\mathcal{L}^{\mathrm{c}}(\mathrm{E})\) 在 \(\mathcal{L}(\mathrm{E})\) 中不是闭的。
\end{enumerate}

\begin{enumerate}
\def\labelenumi{\arabic{enumi})}
\setcounter{enumi}{6}
\item
  设 E 和 F 是局部凸拓扑向量空间。若 F 是可度量化的，则每个紧线性映射 \(E \rightarrow F\) 的像是可数型的。
\item
  构造希尔伯特空间 E 的一个非紧自同态 \(u\) 的例子，使得 \(u^2\) 是紧的。
\item
  对每一对 \((x,y)\in]0,1]^{2}\) ，存在一个整数 \(n\geqslant1\) 和一个整数 k 满足 \(0\leqslant k\leqslant2^{n-1}-1\)，使得
\end{enumerate}

\[
2 ^ {- n} <   x \leqslant 2 ^ {- n + 1}, \quad \frac {2 k}{2 ^ {n}} <   y \leqslant \frac {2 k + 2}{2 ^ {n}}.
\]

然后我们如下定义 N: ]0,1]×]0,1] → C：

\[
\mathrm{N} (x, y) = \left\{ \begin{array}{l l} 2 & \text { if } \frac {2 k}{2 ^ {n}} <   y \leqslant \frac {2 k + 1}{2 ^ {n}} \\ 0 & \text { if } \frac {2 k + 1}{2 ^ {n}} <   y \leqslant \frac {2 k + 2}{2 ^ {n}}. \end{array} \right.
\]

\begin{enumerate}
\def\labelenumi{\alph{enumi})}
\tightlist
\item
  证明公式
\end{enumerate}

\[
u (f) (x) = \int_ {[ 0, 1 ]} \mathrm{N} (x, y) f (y) d y
\]

在 \(L^{1}([0,1])\) 上定义了一个关于勒贝格测度的连续线性映射 u。

\begin{enumerate}
\def\labelenumi{\alph{enumi})}
\setcounter{enumi}{1}
\item
  证明自同态 \(u\) 不是紧的。
\item
  自同态 \(u^2\) 是紧的。
\end{enumerate}

\(d)\) 设 \((x_{n})_{n\in \mathbf{N}}\) 是 \(\mathrm{L}^1 ([0,1])\) 中的一个有界序列。序列 \((u(x_n))_{n\in \mathbf{N}}\) 有一个弱收敛的子列。

\begin{enumerate}
\def\labelenumi{\alph{enumi})}
\setcounter{enumi}{4}
\item
  设 \((x_{n})_{n\in \mathbf{N}}\) 是 \(\mathrm{L}^1 ([0,1])\) 中的一个弱收敛序列。序列 \((u(x_{n}))_{n\in \mathbf{N}}\) 在 \(\mathrm{L}^1 ([0,1])\) 中收敛。
\item
  设 E 是这样的拓扑向量空间：其元素是有界可测的函数 \(f\colon[0,1]\to\mathbf{C}\)，赋以强于 \(\mathrm{L}^{2}([0,1])\) 的拓扑以及相对于 \(\mathrm{L}^{1}([0,1])\) 的弱拓扑的最细拓扑。则由
\end{enumerate}

\[
v (f) (x) = \int_ {[ 0, 1 ]} \overline {{\mathrm{N} (y , x)}} f (y) d y
\]

定义的线性映射 v 是 E 的一个紧自同态。它的转置 \(^{t}v\) 对于 E' 上的强拓扑不是紧的。

\begin{enumerate}
\def\labelenumi{\arabic{enumi})}
\setcounter{enumi}{9}
\tightlist
\item
  设 p 和 r 是满足 \(1 \leqslant p < r\) 的实数。设 u 是从 \(\ell_{\mathrm{K}}^{r}(\mathbf{N})\) 到 \(\ell_{\mathrm{K}}^{p}(\mathbf{N})\) 的一个连续线性映射。
\end{enumerate}

我们用 \(r_{n}\)（相应地，\(p_{n}\)）表示投影

\[
(x _ {i}) _ {i \in \mathbf {N}} \mapsto (x _ {0}, \ldots , x _ {n - 1}, 0, 0, \ldots)
\]

到 \(\ell_{\mathrm{K}}^{r}(\mathbf{N})\) 上（相应地，到 \(\ell_{\mathrm{K}}^{p}(\mathbf{N})\) 上）。我们用 \(\| \cdot \| _p\)（相应地，\(\| \cdot \| _r\)）表示 \(\ell_{\mathrm{K}}^{p}(\mathbf{N})\)（相应地，\(\ell_{\mathrm{K}}^{r}(\mathbf{N})\)）上的范数。

\begin{enumerate}
\def\labelenumi{\alph{enumi})}
\item
  证明对每个 \(n \in \mathbf{N}\) ，有 \(\|p_n u(1 - r_m)\| \to 0\)（当 \(m \to +\infty\)），其中 1 表示 \(\ell_{\mathrm{K}}^{r}(\mathbf{N})\) 的恒同态射。
\item
  对 \(n\) 、 \(m \in \mathbf{N}\) ，令 \(u_{n,m} = (1 - p_n)u(1 - r_m)\) 。设 \(\alpha = (\alpha_n)_{n \in \mathbf{N}}\) 是一个严格正实数的序列，满足 \(\|\alpha\|_r = 1\) 且 \(\alpha\) 不属于 \(\ell_{\mathrm{K}}^p(\mathbf{N})\)，并设 \((\varepsilon_n)\) 是一个严格正实数的序列，使得级数 \(\sum \varepsilon_n\alpha_n\) 收敛。
\end{enumerate}

若 \(u\) 不是紧的，证明存在 \(\delta > 0\) 、整数的严格递增序列 \((n_i)_{i \in \mathbf{N}}\) 与 \((m_i)_{i \in \mathbf{N}}\)，以及两个序列：\((x_i)_{i \in \mathbf{N}}\)（由 \(\ell_{\mathrm{K}}^r(\mathbf{N})\) 的元素组成）与 \((y_i)_{i \in \mathbf{N}}\)（由 \(\ell_{\mathrm{K}}^p(\mathbf{N})\) 的元素组成），具有下列性质：

\begin{enumerate}
\def\labelenumi{(\roman{enumi})}
\item
  有 \(\|x_{i}\|_{r}=1\) 对所有 \(i\in N\) 成立；
\item
  集合 \(S_{i} = \{n \in N \mid x_{i,n} \neq 0\}\) 两两不相交，且集合 \(T_{i} = \{n \in N \mid y_{i,n} \neq 0\}\) 两两不相交；
\item
  有 \(\| y_i\| _p\geqslant \delta\)
\item
  有 \(\|u(x_{i})\|_{q} > \delta - 2\varepsilon_{i}\) 。
\end{enumerate}

（用归纳法进行证明，并利用公式

\[
u = u _ {n, m} + p _ {n} u (1 - r _ {m}) + (1 - p _ {n}) u r _ {m} + p _ {n} u q _ {m}
\]

对 N 中的 n 和 m。）

\begin{enumerate}
\def\labelenumi{\alph{enumi})}
\setcounter{enumi}{2}
\item
  用反证法推出 u 是紧的（"皮特定理"）。
\item
  证明 \(\ell_{\mathrm{K}}^{p}(\mathbf{N})\) 到 \(\ell_{\mathrm{K}}^{r}(\mathbf{N})\) 的典范单射不是紧的。
\item
  设 p 和 r 是满足 \(1 \leqslant p \leqslant r \leqslant +\infty\) 的实数。证明 \(\ell_{\mathrm{K}}^{r}(\mathbf{N})\) 同构于 \(\ell_{\mathrm{K}}^{p}(\mathbf{N})\) 当且仅当 p = r。
\end{enumerate}

\begin{enumerate}
\def\labelenumi{\arabic{enumi})}
\setcounter{enumi}{10}
\tightlist
\item
  设 E 和 F 是 K 上的巴拿赫空间。称线性映射 \(u \in \mathcal{L}(\mathrm{E};\mathrm{F})\) 为完全连续的，如果对 E 中每个弱收敛序列 \((x_{n})_{n \in \mathbb{N}}\)，序列 \((u(x_{n}))_{n \in \mathbb{N}}\) 在 F 中收敛。
\end{enumerate}

\begin{enumerate}
\def\labelenumi{\alph{enumi})}
\item
  每个紧线性映射都是完全连续的。
\item
  若 E 的对偶空间 E' 是可数生成的，则每个完全连续线性映射 E → F 都是紧的。
\item
  设 \((u_{n})_{n\in\mathbf{N}}\) 是 \(\ell_{\mathrm{K}}^{1}(\mathbf{N})\) 中元素的一个序列。证明 \((u_{n})\) 收敛当且仅当它弱收敛（"舒尔定理"；参见 EVT, IV, p. 54, 习题 15）。
\end{enumerate}

\(d)\) 由此推出 \(\ell_{\mathrm{K}}^{1}(\mathbf{N})\) 到 \(\ell_{\mathrm{K}}^{2}(\mathbf{N})\) 的单射是完全连续的，但不是紧的。

\begin{enumerate}
\def\labelenumi{\alph{enumi})}
\setcounter{enumi}{4}
\item
  对每个巴拿赫空间，每个连续线性映射 \(\mathrm{E} \to \ell_{\mathrm{K}}^{1}(\mathbf{N})\) 都是完全连续的，且每个连续线性映射 \(\ell_{\mathrm{K}}^{1}(\mathbf{N}) \to \mathrm{E}\) 都是完全连续的。
\item
  若 E 是自反的或其对偶是可数生成的，则每个连续线性映射 \(\mathrm{E} \rightarrow \ell_{\mathrm{K}}^{1}(\mathbf{N})\) 都是紧的。
\end{enumerate}

\begin{enumerate}
\def\labelenumi{\arabic{enumi})}
\setcounter{enumi}{11}
\tightlist
\item
  设 E 是一个无限维自反巴拿赫空间。设 \(u \in \mathcal{L}^{c}(E)\) 是一个紧自同态。假设 u 不是有限秩的。设 \(E_{\sigma}\) 是赋以弱拓扑的空间 E。
\end{enumerate}

证明线性映射 \(^{t}u\) 是 \((\mathrm{E}_{\sigma})'\) 的一个紧自同态。

\begin{enumerate}
\def\labelenumi{\arabic{enumi})}
\setcounter{enumi}{12}
\tightlist
\item
  设 E 是 K 上的一个可数生成的巴拿赫空间。回顾（EVT, IV, p. 70, 习题 14）：E 中元素的一个序列 \((x_{n})_{n\in \mathbf{N}}\) 称为 E 的绍德尔基，如果对每个 \(x\in \mathrm{E}\) ，存在唯一的纯量序列 \((\lambda_n(x))_{n\in \mathbf{N}}\)，使得
\end{enumerate}

\[
x = \sum_ {n = 1} ^ {+ \infty} \lambda_ {n} (x) x _ {n}
\]

其中级数在 E 中收敛。

\begin{enumerate}
\def\labelenumi{\alph{enumi})}
\item
  每个可分希尔伯特空间都具有巴拿赫基；每个空间 \(\ell^{p}(\mathrm{I})\)（其中 I 可数）都具有巴拿赫基。
\item
  由复平面中半径为 1 的开圆盘 D 上有界全纯函数组成的巴拿赫空间 E，赋以范数 \(\| f \| = \sup_{z \in \mathrm{D}} |f(z)|\)，具有绍德尔基。
\item
  每个具有巴拿赫基的巴拿赫空间 E 都是一个逼近空间。更精确地说，存在实数 \(c \geqslant 0\)，使得 \(1_{E}\) 属于 \(\mathcal{L}_{\mathrm{pc}}(\mathrm{E})\) 中由 E 中范数 \(\leqslant c\) 的有限秩自同态组成的集合的闭包。
\end{enumerate}

\begin{enumerate}
\def\labelenumi{\arabic{enumi})}
\setcounter{enumi}{13}
\tightlist
\item
  \begin{enumerate}
  \def\labelenumii{\alph{enumii})}
  \tightlist
  \item
    设 E 和 F 是希尔伯特空间。从 E 到 F 的连续线性映射 u 是紧的，当且仅当 u 的像不含无限维的闭子空间。
  \end{enumerate}
\end{enumerate}

\begin{enumerate}
\def\labelenumi{\alph{enumi})}
\setcounter{enumi}{1}
\tightlist
\item
  a) 中关于从 E 到 F 的紧自同态的刻画并非对每一对巴拿赫空间 (E, F) 都成立。（可利用 p. 121 的习题 5。）
\end{enumerate}

\begin{enumerate}
\def\labelenumi{\arabic{enumi})}
\setcounter{enumi}{14}
\tightlist
\item
  设 \(E_{0}\) 是由函数 \(f: N^{*} \times N^{*} \to C\) 组成的复向量空间。对 \(k \in N^{*}\) ，设 \(\alpha_{k} \in E_{0}\) 是满足下式的函数：
\end{enumerate}

\[
\alpha_ {k} (n, m) = \left\{ \begin{array}{l l} m ^ {k} & \text { if } 1 \leqslant n \leqslant k - 1 \\ n ^ {k} & \text { if } n \geqslant k. \end{array} \right.
\]

设 E 是 \(E_{0}\) 的由满足 \(p_{k}(f)\) 对所有 \(k \in N^{*}\) 有限的函数 f 组成的向量子空间，其中 \(p_{k}\) 由

\[
p _ {k} (f) = \sum_ {n, m} \alpha_ {k} (n, m) | f (n, m) |.
\]

\begin{enumerate}
\def\labelenumi{\alph{enumi})}
\tightlist
\item
  赋以由范数 \((p_{k})_{k\in\mathbf{N}^{*}}\) 定义的拓扑的空间 E 是一个弗雷歇空间，也是一个蒙泰尔空间。它的对偶 E' 等同于 \(E_{0}\) 中由满足下式的函数 g 组成的子空间：对至少一个 \(k\geqslant1\) ，有
\end{enumerate}

\[
\sup _ {m, n} \alpha_ {k} (m, n) ^ {- 1} | g (m, n) | <   + \infty
\]

（参见 EVT, IV, p. 63, 习题 8）。

\begin{enumerate}
\def\labelenumi{\alph{enumi})}
\setcounter{enumi}{1}
\tightlist
\item
  考虑由下式定义的线性映射 \(u\)，它把 E 映到 \(\ell^1 (\mathbf{N}^*)\)：
\end{enumerate}

\[
u (f) = \left(\sum_ {n} f (n, m)\right) _ {m \in \mathbf {N} ^ {*}}.
\]

线性映射 u 是连续的。

\begin{enumerate}
\def\labelenumi{\alph{enumi})}
\setcounter{enumi}{2}
\item
  把 \(\ell^{\infty}(\mathbf{N}^{*})\) 等同于 \(\ell^{1}(\mathbf{N}^{*})\) 的对偶。\(u\) 的转置是映射 \(v\colon \ell^{\infty}(\mathbf{N}^{*})\to \mathrm{E}^{\prime}\subset \mathrm{E}_{0}\)，由 \((x_{n})\mapsto f\) 确定，其中 \(f(n,m) = x_{m}\) 对所有 \((n,m)\in \mathbf{N}^{*}\times \mathbf{N}^{*}\) 成立。
\item
  映射 u 诱导一个同构 \(\mathrm{E}/\mathrm{Ker}(u)\to\ell^{1}(\mathbf{N}^{*})\)（参见出处同上）。
\item
  映射 u 不是紧的。
\item
  u 的转置 v 对于在 \(\ell^{\infty}(\mathbf{N}^{*})\) 的有界子集上收敛的拓扑是紧的。
\item
  映射 v 是分离局部凸空间之间一个紧线性映射的例子，其转置不是紧的。
\end{enumerate}

\begin{enumerate}
\def\labelenumi{\arabic{enumi})}
\setcounter{enumi}{15}
\item
  设 E 是局部凸空间，u 是 E 的一个紧自同态。用 \(E_{b}^{\prime}\) 表示赋有有界收敛拓扑的 E 的对偶空间。线性映射 \(^{t}u \circ ^{t}u\) 是 \(E_{b}^{\prime}\) 的一个紧自同态。
\item
  设 E 是可数型的巴拿赫空间。则 \(\mathcal{L}^{\mathrm{c}}(\mathrm{E})\) 赋以由 \(\mathcal{L}(\mathrm{E})\) 的范数诱导的拓扑后是可数型的，当且仅当对偶空间 \(\mathrm{E}'\) 是可数型的。
\item
  设 E、F、G 是巴拿赫空间，u 是 \(\mathcal{L}^{c}(\mathrm{E};\mathrm{F})\) 的一个元素。
\end{enumerate}

\begin{enumerate}
\def\labelenumi{\alph{enumi})}
\tightlist
\item
  设 v 是 \(\mathcal{L}(\mathrm{E};\mathrm{G})\) 的一个元素。假设对每个 \(\varepsilon > 0\) ，存在 C \textgreater{} 0，使得
\end{enumerate}

\[
\| v (x) \| _ {\mathrm{G}} \leqslant \varepsilon \| x \| _ {\mathrm{E}} + \mathrm{C} \| u (x) \| _ {\mathrm{F}}
\]

对所有 \(x\in \mathrm{E}\) 成立。则 \(v\in \mathcal{L}^{\mathrm{c}}(\mathrm{E};\mathrm{G})\) 。

\begin{enumerate}
\def\labelenumi{\alph{enumi})}
\setcounter{enumi}{1}
\tightlist
\item
  设 w 是 \(\mathcal{L}(\mathrm{F};\mathrm{G})\) 的一个单射元素。则对每个 \(\varepsilon > 0\) ，存在 C \textgreater{} 0，使得
\end{enumerate}

\[
\| u (x) \| _ {\mathrm{F}} \leqslant \varepsilon \| x \| _ {\mathrm{E}} + \mathrm{C} \| (w \circ u) (x) \| _ {\mathrm{G}}
\]

对所有 \(x \in E\) 成立。

\begin{enumerate}
\def\labelenumi{\arabic{enumi})}
\setcounter{enumi}{18}
\tightlist
\item
  设 X 是一个局部紧、可度量化且在无穷远处可数的拓扑空间，\(\mu\) 是 X 上的一个正测度。我们用 A 表示 \(\mu\) 的原子集。设 \(f \in \mathcal{L}^{\infty}(X, \mu)\) 且 \(p \in [1, +\infty]\) 。
\end{enumerate}

证明乘以 \(f\) 在 \(\mathrm{L}^{p}(\mathrm{X},\mu)\) 中定义一个紧自同态，当且仅当下列两个条件被满足：

\begin{enumerate}
\def\labelenumi{(\roman{enumi})}
\item
  集合 \(\{x\in X-A\mid f(x)\neq0\}\) 是 \(\mu\) -可忽略的；
\item
  对每个 \(\varepsilon > 0\) ，存在 A 的有限子集 F，使得 \(|f(a)| \leqslant \varepsilon\) 对所有 \(a \in \mathbf{A} - \mathbf{F}\) 成立。
\end{enumerate}

（考虑乘以 \(f\) 的运算在空间 \(\mathrm{L}^{p}(\mathrm{X}_{\varepsilon},\mu)\) 上的作用，其中 \(\mathrm{X}_{\varepsilon}\) 是由 \(x\in\mathrm{X}\) 且满足 \(|f(x)|\geqslant\varepsilon\) 的元素组成的集合。）

\begin{enumerate}
\def\labelenumi{\arabic{enumi})}
\setcounter{enumi}{19}
\tightlist
\item
  设 X、Y 是紧度量空间，T 是从 \(\mathcal{C}(\mathrm{Y})\) 到 \(\mathcal{C}(\mathrm{X})\) 的一个紧算子。
\end{enumerate}

证明存在一个正测度 \(\mu\)（在 \(X\) 上）以及一个可测映射 \(N\)（从 \(X \times Y\) 到 \(\mathbf{C}\)），使得

\begin{enumerate}
\def\labelenumi{\alph{enumi})}
\item
  对每个 \(y \in Y\) ，映射 \(\mathrm{N}_{y} \colon x \mapsto \mathrm{N}(x, y)\) 是从 \(X\) 到 \(\mathbf{C}\) 的 \(\mu\) -可积映射；
\item
  映射 \(y \mapsto N_{y}\) 是从 Y 到 \(\mathcal{L}^{1}(X, \mu)\) 的连续映射；
\item
  对每个 \(f \in \mathcal{C}(\mathrm{X})\) 和每个 \(y \in \mathrm{Y}\) ，我们有
\end{enumerate}

\[
\mathrm{T} f (y) = \int_ {\mathrm{X}} \mathrm{N} (x, y) f (x) d \mu (x).
\]

\begin{enumerate}
\def\labelenumi{\arabic{enumi})}
\setcounter{enumi}{20}
\tightlist
\item
  我们用 D 表示 C 中以 0 为心、以 1 为半径的开圆盘，用 H \(^{2}\) (D) 表示 D 上的哈代空间（EVT, V, p. 7, 例 5）。我们记 T = R/2πZ，并赋以规范化的哈尔测度。设 j 是由下式定义的从 H \(^{2}\) (D) 到 L \(^{2}\) (T) 的等距：
\end{enumerate}

\[
j (f) (\theta) = \sum_ {n \geqslant 0} \frac {f ^ {(n)} (0)}{n !} e ^ {i n \theta}.
\]

H \(^{2}\) (D) 在 \(j\) 下的像将记为 H \(^{2}\) (T)。

设 \(\widetilde{\mathrm{P}}\) 是 \(\mathrm{L}^2 (\mathbf{T})\) 的以 \(\mathrm{H}^2 (\mathbf{T})\) 为像的正交投影，并设 P 是线性映射 \(f\mapsto \widetilde{\mathrm{P}} (f)\)，它从 \(\mathrm{L}^2 (\mathbf{T})\) 到 \(\mathrm{H}^2 (\mathbf{T})\)。若 \(f\) 是 \(\mathrm{L}^{\infty}(\mathbf{T})\) 的一个元素，设 \(\mathrm{M}_f\) 是自同态 \(g\mapsto fg\)，它作用于 \(\mathrm{L}^2 (\mathbf{T})\)。

证明若 \(f\) 是连续函数，则换位子 \(\widetilde{\mathrm{P}}\mathrm{M}_f - \mathrm{M}_f\widetilde{\mathrm{P}}\) 是一个紧线性映射。（考虑 \(f\) 是三角多项式的特殊情形。）

\begin{enumerate}
\def\labelenumi{\arabic{enumi})}
\setcounter{enumi}{21}
\tightlist
\item
  我们保留上一习题的记号。设 R 是从 \(\mathrm{H}^{2}(\mathbf{T})\) 到 \(\mathrm{L}^{2}(\mathbf{T})\) 的由 \(\mathrm{R}f(\theta)=f(-\theta)\) 定义的连续线性映射。对 \(g\in\mathrm{L}^{\infty}(\mathbf{T})\)，我们令 \(\Gamma_{g}=\mathrm{PM}_{g}\mathrm{R}\) 。
\end{enumerate}

\begin{enumerate}
\def\labelenumi{\alph{enumi})}
\item
  证明线性映射 \(\Gamma\colon g\mapsto\Gamma_{g}\) 是从 \(\mathrm{L}^{\infty}(\mathbf{T})\) 到 \(\mathcal{L}(\mathrm{H}^{2}(\mathbf{T}))\) 的连续映射，范数 \(\leqslant 1\)，且其核 F 等于 \(\mathrm{H}^{2}(\mathbf{T})^{\circ}\cap\mathrm{L}^{\infty}(\mathbf{T})\) 。
\item
  若 \(g\) 是连续函数，则 \(\Gamma_g\) 是 \(\mathrm{H}^2 (\mathbf{T})\) 的一个紧自同态。
\item
  设 \(g \in \mathrm{L}^{\infty}(\mathbf{T})\) 。我们用 \(\alpha_{g}\) 表示 \(\mathrm{H}^{2}(\mathbf{T})\) 上由下式定义的线性形式
\end{enumerate}

\[
\alpha_ {g} (f) = \int_ {\mathbf {T}} \Gamma_ {g} (f) (\theta) d \theta .
\]

设 \(\mathcal{A}\) 是在 \(j\) 下 \(\mathrm{H}^2 (\mathrm{D})\) 中由可解析延拓到 \(\overline{\mathrm{D}}\) 的一个邻域的全纯函数组成的子空间的像。证明若 \(f_{1}\) 和 \(f_{2}\) 属于 \(\mathcal{A}\) ，则

\[
\alpha_ {g} (f _ {1} f _ {2}) = \int_ {\mathbf {T}} \Gamma_ {g} (f _ {1}) (\theta) f _ {2} (- \theta) d \theta .
\]

\(d)\) 设 \(\varphi\) 是 \(\overline{\mathrm{D}}\) 的一个邻域上的全纯函数，在 \(\partial\mathrm{D}\) 上不为零，并设 \(f = j(\varphi|\mathrm{D})\) 。则存在 \(f_{1}\) 和 \(f_{2}\)（都在 \(\mathscr{A}\) 中），使得 \(f = f_{1}f_{2}\)，且

\[
\| f \| _ {\mathrm{L} ^ {1} (\mathbf {T})} = \| f _ {1} \| _ {\mathrm{L} ^ {2} (\mathbf {T})} \| f _ {2} \| _ {\mathrm{L} ^ {2} (\mathbf {T})}.
\]

（设 \(z_{1},\ldots,z_{n}\) 是 \(\varphi\) 在 D 中的零点；则存在一个在 D 的邻域内全纯的函数 \(\psi\)，使得

\[
\varphi (z) = \prod_ {k = 1} ^ {n} \frac {z - z _ {k}}{1 - \overline {{z}} _ {k} z} \psi (z) ^ {2}
\]

对所有 z 成立。）

\begin{enumerate}
\def\labelenumi{\alph{enumi})}
\setcounter{enumi}{4}
\tightlist
\item
  证明对所有 \(f \in \mathrm{H}^{2}(\mathbf{T})\) ，我们有不等式
\end{enumerate}

\[
\left| \alpha_ {g} (f) \right| \leqslant \left\| \Gamma_ {g} \right\| _ {\mathscr {L} \left(\mathrm{H} ^ {2} (\mathbf {T})\right)} \| f \| _ {\mathrm{L} ^ {1} (\mathbf {T})}.
\]

（若 \(\varphi \in \mathrm{H}^2 (\mathrm{D})\) ，则存在 ]0,1[ 中的一个序列 \((r_n)_{n\in \mathbf{N}}\)，使得对每个 \(n\in \mathbf{N}\) ，函数 \(\varphi_{n}\)（由 \(\varphi_{n}(z) = \varphi (r_{n}z)\) 定义）在 \(\partial \mathrm{D}\) 上没有零点，且 \(\lim \varphi_n = \varphi\) 在 \(\mathrm{H}^2 (\mathrm{D}).\) 中

由此推出存在 \(h \in \mathrm{L}^{\infty}(\mathbf{T})\)，使得

\[
\alpha_ {g} (f) = \int_ {\mathbf {T}} f (\theta) h (\theta) d \theta
\]

对所有 \(f \in \mathrm{H}^2(\mathbf{T})\) 成立，且 \(\|h\|_{\mathrm{L}^\infty(\mathbf{T})} \leqslant \| \Gamma_g \|_{\mathcal{L}(\mathrm{H}^2(\mathbf{T}))}\) 。

\(f)\) 证明存在 \(k \in \mathrm{L}^{\infty}(\mathbf{T})\)，使得 \(\Gamma_g = \Gamma_k\) 且 \(\| k \|_{\mathrm{L}^{\infty}(\mathbf{T})} = \| \Gamma_g \|_{\mathcal{L}(\mathrm{H}^2(\mathbf{T}))}\) 。

于是映射 \(\widehat{\Gamma}\colon\mathrm{L}^{\infty}(\mathbf{T})/\mathrm{F}\to\mathcal{L}(\mathrm{H}^{2}(\mathbf{T}))\) 是一个等距（“内哈里定理”）。

\(g)\) 设 \(g\) 是 \(\mathbf{T}\) 上的连续函数。则我们有 \(d(g, \mathrm{F}) = d(g, \mathrm{F} \cap \mathcal{C}(\mathbf{T}))\) 。（对 \(\mathrm{R} > 1\) ，用下式定义线性映射 \(\mathrm{P}_{\mathrm{R}}\)，即从 \(\mathrm{L}^{\infty}(\mathbf{T})\) 到 \(\mathcal{C}(\mathbf{T})\) 的映射：

\[
\mathrm{P} _ {\mathrm{R}} f (\theta) = \int_ {\mathbf {T}} \frac {\mathrm{R} ^ {2} - 1}{| \mathrm{R} e ^ {i (\theta - \varphi)} - 1 | ^ {2}} f (\varphi) d \varphi .
\]

则有 \(\| \mathrm{P}_{\mathrm{R}}f\|_{\mathcal{C}(\mathbf{T})}\leqslant \| f\|_{\mathrm{L}^{\infty}(\mathbf{T})}\)，且若 \(f\in \mathcal{C}(\mathbf{T})\) ，则 \(\mathrm{P_R}f\to f\) 在 \(\mathcal{C}(\mathbf{T})\) 中成立（当 \(\mathrm{R}\rightarrow 1.\)）。

由此推出 \(\Gamma(\mathcal{C}(\mathbf{T}))\) 是 \(\mathcal{L}(\mathrm{H}^{2}(\mathbf{T}))\) 的一个闭子空间（“萨拉森定理”）。

\begin{enumerate}
\def\labelenumi{\alph{enumi})}
\setcounter{enumi}{7}
\tightlist
\item
  设 \(g \in \mathrm{L}^{\infty}(\mathbf{T})\) 满足 \(\Gamma_{g}\) 是一个紧自同态。对 \(m \in Z\)，我们令 \(e_{m}(\theta) = e^{im\theta}\) 。证明 \(\|\Gamma_{ge_{-m}}\|_{\mathcal{L}(\mathrm{H}^{2}(\mathbf{T}))}\) 当 m 趋于 \(+\infty\) 时趋于零，然后证明存在 F 中的一个序列 \((h_{m})_{m \in \mathbf{N}}\)，使得序列 \((e_{m}h_{m})_{m \in \mathbf{N}}\) 在 \(\mathrm{L}^{\infty}(\mathbf{T})\) 中收敛于 g。由此推出在 g 模 F 的类中存在一个连续函数（“哈特曼定理”）。
\end{enumerate}

\begin{enumerate}
\def\labelenumi{\arabic{enumi})}
\setcounter{enumi}{22}
\item
  设 B 是一个无限维巴拿赫空间。我们用 E 表示 B 的强对偶，用 F 表示赋有拓扑 \(\sigma(\mathrm{B}^{\prime},\mathrm{B})\) 的 B 的对偶。证明 B 的对偶的恒同映射从 E 到 F 是紧的，但其转置从 \(\mathrm{F}_{b}^{\prime}\) 到 \(\mathrm{E}_{b}^{\prime}\) 不是紧的。
\item
  设 E 是一个实或复的巴拿赫空间，\(u \in \mathcal{L}(\mathrm{E})\) 。假设序列 \((u^{n})_{n \in \mathbf{N}}\) 有界，且 u 从 E 到赋有弱拓扑 \(\sigma(\mathrm{E}, \mathrm{E}')\) 的 E 是紧的。对每个整数 \(N \geqslant 1\)，我们令
\end{enumerate}

\[
v _ {\mathrm{N}} = \frac {1}{\mathrm{N}} (1 _ {\mathrm{E}} + u + \dots + u ^ {\mathrm{N-1}}) \in \mathcal {L} (\mathrm{E}).
\]

\begin{enumerate}
\def\labelenumi{\alph{enumi})}
\item
  设 F 是 \(1_{\mathrm{E}} - u\) 的像的闭包。对每个 \(x \in \mathrm{F}\) ，我们有 \(v_{\mathrm{N}}(x) \to 0\)（在 E 中，当 \(\mathrm{N} \to +\infty\) 时）。
\item
  设 \(x \in E\) 。序列 \((v_{\mathrm{N}}(x))\) 收敛到一个元素 \(y \in E\)，使得 \(u(y) = y\) 。（存在 \((v_{\mathrm{N}}(x))\) 的一个子序列，它弱收敛到 \(y \in E\) ；利用哈恩–巴拿赫定理证明 \(x - y \in F\) 。）
\item
  映射 \(w \colon \mathrm{E} \to \mathrm{E}\) 满足 \(w(x)\) 是上一问题给出的元素 \(y\)，它是一个连续投影，其像是 \(u\) 的对应于特征值 1 的特征空间。
\end{enumerate}

\(d)\) 设 \(z \in \mathbf{C}\) 满足 \(|z| = 1\) 。由下式定义的 \(\mathcal{L}(\mathrm{E})\) 中的序列

\[
v _ {z, \mathrm{N}} = \frac {1}{\mathrm{N}} (1 _ {\mathrm{E}} + z ^ {- 1} u \dots + z ^ {- (\mathrm{N-1})} u ^ {\mathrm{N-1}})
\]

在赋有逐点收敛拓扑的 \(\mathcal{L}(\mathrm{E})\) 中收敛。用 \(w_{z}\) 表示它的极限。则 \(w_{z}\) 是一个投影，我们有 \(w_{z}w = ww_{z} = zw_{z}\)，以及 \(w_{z_1}w_{z_2} = 0\)（若 \(z_{1} \neq z_{2}\)）。

\begin{enumerate}
\def\labelenumi{\arabic{enumi})}
\setcounter{enumi}{24}
\tightlist
\item
  \begin{enumerate}
  \def\labelenumii{\alph{enumii})}
  \tightlist
  \item
    设 \(N \geqslant 1\) 是一个整数。我们用 \(G_{N}\)（相应地，\(H_{N}\)）表示集合 \(\{-1,1\}^{N}\)（相应地，集合 \(\{-1,2\}^{N}\)）。我们用 \(\mu_{N}\)（相应地，\(\nu_{N}\)）表示 \(G_{N}\) 上（相应地，\(H_{N}\) 上）总质量为 1 的正测度，使得对每个 \((t_{i})_{1 \leqslant i \leqslant N}\)（属于 \(G_{N}\)，相应地，属于 \(H_{N}\)），有 \(\mu_{\mathrm{N}}((t_{i})_{1 \leqslant i \leqslant N}) = \frac{1}{2^{\mathrm{N}}}\) ，相应地，
  \end{enumerate}
\end{enumerate}

\[
\nu_ {\mathrm{N}} ((t _ {i}) _ {1 \leqslant i \leqslant \mathrm{N}}) = \left(\frac {1}{3}\right) ^ {j} \left(\frac {2}{3}\right) ^ {\mathrm{N} - j},
\]

其中 j 是使得 \(t_{i}=2\) 的指标 i 的个数。

存在一个实数 \(c > 0\)，使得对每个由复数组成的族 \((a_{i})_{1\leqslant i\leqslant N}\)，我们有

\[
\mu_ {\mathrm{N}} \Big (\Big \{(t _ {i}) _ {1 \leqslant i \leqslant \mathrm{N}} \in \mathrm{G} _ {\mathrm{N}} | \Big | \sum_ {i = 1} ^ {\mathrm{N}} a _ {i} t _ {i} \Big | > c (\log \mathrm{N}) ^ {1 / 2} \Big (\sum_ {i = 1} ^ {\mathrm{N}} | a _ {i} | ^ {2} \Big) ^ {1 / 2} \Big \} \Big) <   \frac {c}{\mathrm{N} ^ {3}},
\]

以及

\[
\nu_ {\mathrm{N}} \Big (\Big \{(t _ {i}) _ {1 \leqslant i \leqslant \mathrm{N}} \in \mathrm{H} _ {\mathrm{N}} | \left| \sum_ {i = 1} ^ {\mathrm{N}} a _ {i} t _ {i} \right| > c (\log \mathrm{N}) ^ {1 / 2} \Big (\sum_ {i = 1} ^ {\mathrm{N}} | a _ {i} | ^ {2} \Big) ^ {1 / 2} \Big \} \Big) <   \frac {c}{\mathrm{N} ^ {3}}.
\]

（我们可以假设 \(a_{i} \in \mathbf{R}\) 对所有 \(i\) 成立，且 \(\sum |a_{i}|^{2} = 1\) 。对第一个不等式，证明并利用上界估计

\[
\int_ {G _ {N}} \exp \left(\lambda \left| \sum_ {i = 1} ^ {N} a _ {i} t _ {i} \right|\right) d \mu_ {N} \leqslant 2 \exp (\lambda^ {2})
\]

对每个实数 \(\lambda > 0\) 成立；对第二个上界估计，利用不等式 \(\frac{1}{3}(e^{2x} + 2e^{-x}) \leqslant e^{2x^2}\)（对所有 \(x \in \mathbf{R}\)）以类似的方式推理。）

\begin{enumerate}
\def\labelenumi{\alph{enumi})}
\setcounter{enumi}{1}
\tightlist
\item
  对 \(k \in N\)，令 \(\mathrm{A}_{k} = \mathbf{Z}/(3 \cdot 2^{k})\mathbf{Z}\) 。存在一个不依赖于 k 的实数 \(c' > 0\)，以及一个函数 \(f_{k} \colon \widehat{\mathrm{A}}_{k} \to \{-1, 2\}\)，使得
\end{enumerate}

\[
\sum_ {\chi \in \widehat {\mathrm{A}} _ {k}} f _ {k} (\chi) = 0
\]

且对所有 \(x \in A_{k}\) ，有

\[
\Big | \sum_ {\chi \in \widehat {\mathrm{A}} _ {k}} \chi (x) f _ {k} (\chi) \Big | \leqslant c ^ {\prime} (k + 1) ^ {1 / 2} 2 ^ {k / 2}.
\]

（利用 \(a\)），确定满足第二个条件的函数 \(f_{k}\) 的存在性，然后修改 \(f_{k}\) 以得到第一个条件。）

我们用 \(B_{k}\)（相应地，\(C_{k}\)）表示特征标 \(\chi\)（\(A_{k}\) 上的）中满足 \(f_{k}(\chi)=2\)（相应地，\(f_{k}(\chi)=-1\)）者所成的集合。我们有 \(\mathrm{Card(B_{k})}=2^{k}\) 和 \(\mathrm{Card(C_{k})}=2^{k+1}\) 。

对每个 \(k \geqslant 1\) ，我们固定一个双射 \(\iota_k \colon B_k \to C_{k-1}\) 。

我们用 A 表示诸 \(A_{k}\)（\(k \in N\)）的不交并；我们赋给 A 离散拓扑，并记 \(\mathrm{E} = \mathcal{C}_{b}(\mathrm{A})\) 。

我们固定一个函数族 \(\varepsilon = (\varepsilon_k)_{k\geqslant \mathbf{N}}\)，其中的函数 \(\varepsilon_{k}\colon \mathrm{B}_{k}\to \{-1,1\}\)。然后我们定义函数 \(e_{k,\chi}\in \mathrm{E}\)（对 \(k\in \mathbf{N}\) 和 \(\chi \in \mathrm{B}_k\)）如下：

\[
e _ {k, \chi} (x) = \left\{ \begin{array}{l l} \iota_ {k} (\chi) (x) & \text { if } k \geqslant 1 \text {   and   } x \in \mathrm{A} _ {k - 1}, \\ \varepsilon_ {k} (\chi) \chi (x) & \text { if } x \in \mathrm{A} _ {k} \\ 0 & \text { otherwise. } \end{array} \right.
\]

我们用 \(E_{\varepsilon}\) 表示由函数 \((e_{k,\chi})_{k\in\mathbf{N},\chi\in\mathbf{B}_{k}}\) 生成的 E 的闭子空间。这是一个可数型的复巴拿赫空间。习题的其余部分旨在证明：对 \(\varepsilon\) 的适当选取，空间 \(\mathrm{E}_{\varepsilon}\) 不具有逼近性质；它不拥有巴拿赫基。

\begin{enumerate}
\def\labelenumi{\alph{enumi})}
\setcounter{enumi}{2}
\tightlist
\item
  对 \(k \in N\) 和 \(\chi \in B_{k}\)，存在一个连续线性形式 \(\lambda_{k,\chi} \in E_{\varepsilon}'\)，使得对每个 \(f \in E_{\varepsilon}\)，
\end{enumerate}

\[
\lambda_ {k, \chi} (f) = \frac {1}{\operatorname{Card} (\mathrm{A} _ {k})} \sum_ {x \in \mathrm{A} _ {k}} \varepsilon_ {k} (\chi) \chi (x ^ {- 1}) f (x).
\]

对 \(k \geqslant 1\) 和 \(f \in \mathrm{E}_{\varepsilon}\) ，我们有

\[
\lambda_ {k, \chi} (f) = \frac {2}{\operatorname{Card} (\mathrm{A} _ {k})} \sum_ {x \in \mathrm{A} _ {k}} \iota_ {k} (\chi) (x ^ {- 1}) f (x).
\]

\begin{enumerate}
\def\labelenumi{\alph{enumi})}
\setcounter{enumi}{3}
\tightlist
\item
  对每个 \(k \in N\) ，用 \(\beta_{k}\) 表示 \(\mathcal{L}(E_{\varepsilon})\) 上由下式定义的连续线性形式
\end{enumerate}

\[
\beta_ {k} (u) = \frac {1}{2 ^ {k}} \sum_ {\chi \in \mathrm{B} _ {k}} \lambda_ {k, \chi} (u (e _ {k, \chi})).
\]

我们有 \(\beta_{k}(1_{\mathrm{E}_{\varepsilon}}) = 1\)

\begin{enumerate}
\def\labelenumi{\alph{enumi})}
\setcounter{enumi}{4}
\tightlist
\item
  设 \(u \in \mathcal{L}(\mathrm{E}_{\varepsilon})\) 且 \(k \in \mathbf{N}\) 。我们有
\end{enumerate}

\[
\beta_ {k + 1} (u) - \beta_ {k} (u) = \frac {1}{\operatorname{Card} \left(\mathrm{A} _ {k}\right)} \sum_ {x \in \mathrm{A} _ {k}} u \left(\varphi_ {k, x}\right) (x)
\]

其中，对 \(x \in \mathrm{A}_k\) ，我们令

\[
\varphi_ {k, x} = \frac {1}{2 ^ {k + 1}} \sum_ {\chi \in B _ {k + 1}} \iota_ {k + 1} (\chi) (x ^ {- 1}) e _ {k + 1, \chi} - \frac {1}{2 ^ {k}} \sum_ {\chi \in B _ {k}} \varepsilon_ {k} (\chi) \chi (x ^ {- 1}) e _ {k, \chi}.
\]

我们有 \(|\varphi_{k,x}(y)| \leqslant c'(k + 1)^{1/2}2^{-k/2}\)，对 \((x,y) \in \mathrm{A}_k^2\) 成立。

\begin{enumerate}
\def\labelenumi{\alph{enumi})}
\setcounter{enumi}{5}
\item
  存在一个实数 \(c'' \geqslant c'\)，使得对每个 \(k \geqslant 1\) ，存在一个函数 \(\varepsilon_k \colon B_k \to \{-1,1\}\)，满足 \(|\varphi_{k,x}(y)| \leqslant c''3(k+1)^{1/2}2^{-k/2}\) 对 \((x,y) \in A_k \times A_{k-1}\) 成立。（应用 a) 的第一个不等式。）我们现在假设 \(\varepsilon = (\varepsilon_k)\) 满足这些性质。
\item
  对每个 \(k \in \mathbf{N}\) 和 \(x \in \mathrm{A}_k\) ，我们有 \(\| \varphi_{k,x} \| \leqslant c''(k + 1)^{1/2} 2^{-k/2}\) 。
\item
  设 K 是 \(E_{\varepsilon}\) 的子集，其元素为 \(e_{0,\chi_{0}}\)（其中 \(\chi_{0}\) 是 \(B_{0}\) 的唯一元素）以及函数 \((k+1)^{2}\varphi_{k,x}\)（\(k\in N\)，\(x\in A_{k}\)）。集合 K 在 \(E_{\varepsilon}\) 中是相对紧的，且我们有
\end{enumerate}

\[
| \beta_ {k + 1} (u) - \beta_ {k} (u) | \leqslant \frac {1}{(k + 1) ^ {2}} \sup _ {x \in K} \| u (x) \|
\]

对所有 \(k \in \mathbf{N}\) 和每个 \(u \in \mathcal{L}(\mathrm{E}_{\varepsilon})\) 成立。

\begin{enumerate}
\def\labelenumi{\roman{enumi})}
\tightlist
\item
  序列 \((\beta_{k})_{k\in \mathbf{N}}\) 在赋有逐点收敛拓扑的空间 \(\mathcal{L}(\mathrm{E}_{\varepsilon})\) 的对偶中收敛。用 \(\beta\) 表示它的极限。对每个 \(u\in \mathcal{L}(\mathrm{E}_{\varepsilon})\) ，我们有 \(|\beta (u)|\leqslant 3\sup_{x\in \mathrm{K}}\| u(x)\|\) 。
\end{enumerate}

\begin{enumerate}
\def\labelenumi{\alph{enumi})}
\setcounter{enumi}{9}
\tightlist
\item
  我们有 \(\beta(1_{\mathrm{E}_{\varepsilon}})=1\)，且 \(\beta(u)=0\)（若 \(u\) 是有限秩的）。
\end{enumerate}

\(k)\) 空间 \(\mathrm{E}_{\varepsilon}\) 不具有逼近性质。

\begin{enumerate}
\def\labelenumi{\alph{enumi})}
\setcounter{enumi}{11}
\tightlist
\item
  设 \(p \in]2, +\infty[\) 。我们赋给 A 计数测度，并用 \(E_{\varepsilon}^{p}\) 表示 \(L^{p}(A)\) 的由 \((e_{k,\chi})_{k \in \mathbf{N}, \chi \in \mathrm{B}_{k}}\) 生成的闭子空间。空间 \(E_{\varepsilon}^{p}\) 不具有逼近性质。（验证 \(\varphi_{k,x}\) 的支撑的基数 \(\leqslant 21 \cdot 2^{k}\)，并由此推出 \(\|\varphi_{k,x}\|_{p} = \mathrm{O}((k+1)^{1/2}2^{-k(p-2)/2p})\) 对 \(k \in N\) 和 \(x \in A_{k}\) 成立。）
\end{enumerate}

\begin{enumerate}
\def\labelenumi{\arabic{enumi})}
\setcounter{enumi}{25}
\tightlist
\item
  设 E 和 F 是巴拿赫空间，u 是从 E 到 F 的一个连续线性映射。
\end{enumerate}

\begin{enumerate}
\def\labelenumi{\alph{enumi})}
\tightlist
\item
  证明下列条件等价：
\end{enumerate}

\begin{enumerate}
\def\labelenumi{(\roman{enumi})}
\item
  映射 u 是紧的；
\item
  存在一个序列 \((x_n')_n\)（在 E 的强对偶 \(\mathrm{E}'\) 中趋于 0），满足 \(\| u(x)\| \leqslant \sup |\langle x,x_n' \rangle|\) 对所有 \(x \in \mathrm{E}\) 成立；
\item
  存在巴拿赫空间 \(c_{0}\) 的一个闭子空间 L 以及紧线性映射 \(f: E \to L\) 和 \(g: L \to F\)，使得 \(u = g \circ f\)（若 u 是紧的，对 F 的单位球的极集在 \(^{t}u\) 下的像应用 EVT, IV, p. 24 的推论 1。在假设 (ii) 下，设 \((\lambda_{n})\) 是 \(c_{0}\) 的一个元素，使得序列 \((\lambda_{n}^{-1}x_{n}^{\prime})\) 趋于 0；由 \(f(x) = (\langle x, \lambda_{n}^{-1}x_{n}^{\prime} \rangle)_{n}\) 定义 f，并取 L 为 f 的像的闭包。）
\end{enumerate}

\begin{enumerate}
\def\labelenumi{\alph{enumi})}
\setcounter{enumi}{1}
\tightlist
\item
  若 \(u\) 是紧的，我们有 \(\|u\| = \inf(\sup_n \|x_n'\|)\)，其中下确界取遍满足 (ii) 的序列 \((x_n')\) 所成的集合，且 \(\|u\| = \inf(\|f\|\|g\|)\) ，下确界取遍满足 (iii) 的条件的三元组 (L, \(f,g\) ) 所成的集合。
\end{enumerate}

\begin{enumerate}
\def\labelenumi{\arabic{enumi})}
\setcounter{enumi}{26}
\tightlist
\item
  设 E 和 F 是局部凸空间。映射 \(u \in \mathcal{L}(\mathrm{E};\mathrm{F})\) 称为弱紧的，如果它是从 E 到赋有拓扑 \(\sigma(\mathrm{F},\mathrm{F}')\) 的 F 的紧映射。
\end{enumerate}

\begin{enumerate}
\def\labelenumi{\alph{enumi})}
\item
  弱紧映射构成一个子空间 \(\mathcal{L}^{\mathrm{fc}}(\mathrm{E};\mathrm{F})\)，它是 \(\mathcal{L}(\mathrm{E};\mathrm{F})\) 的子空间。设 \(\mathrm{E}_1\) 、\(\mathrm{F}_1\) 是局部凸空间，\(f\in \mathcal{L}(\mathrm{E}_1;\mathrm{E})\)，\(g\in \mathcal{L}(\mathrm{F};\mathrm{F}_1)\) 且 \(u\in \mathcal{L}(\mathrm{E};\mathrm{F})\)；若 \(u\) 是弱紧的，则 \(g\circ u\circ f\) 亦然。
\item
  我们有 \(\mathcal{L}^{c}(\mathrm{E};\mathrm{F})\subset\mathcal{L}^{\mathrm{fc}}(\mathrm{E};\mathrm{F})\) 。若 E 是无限维自反巴拿赫空间，则 E 的恒同映射是弱紧的但不是紧的。
\end{enumerate}

\begin{enumerate}
\def\labelenumi{\arabic{enumi})}
\setcounter{enumi}{27}
\tightlist
\item
  设 E 和 F 是局部凸空间，且 \(u \in \mathcal{L}(\mathrm{E};\mathrm{F})\) 。
\end{enumerate}

\begin{enumerate}
\def\labelenumi{\alph{enumi})}
\tightlist
\item
  证明下列性质等价：
\end{enumerate}

\begin{enumerate}
\def\labelenumi{(\roman{enumi})}
\item
  E 的每个有界子集在 \(u\) 下的像对于弱拓扑是相对紧的；
\item
  映射 \(^{tt}u\) 把 F'' 映到 E 中。
\end{enumerate}

\begin{enumerate}
\def\labelenumi{\alph{enumi})}
\setcounter{enumi}{1}
\tightlist
\item
  若 F 是拟完备的，证明这些性质等价于：
\end{enumerate}

\begin{enumerate}
\def\labelenumi{(\roman{enumi})}
\setcounter{enumi}{2}
\tightlist
\item
  映射 \(^{t}u\) 把 F' 的每个等度连续子集映到 E' 的一个对于 \(\sigma(\mathrm{E}',\mathrm{E}'')\) 相对紧的子集。
\end{enumerate}

\begin{enumerate}
\def\labelenumi{\alph{enumi})}
\setcounter{enumi}{2}
\tightlist
\item
  若 E 和 F 是巴拿赫空间，则 (i) 至 (iii) 也等价于
\end{enumerate}

(i') 映射 u 是弱紧的；

(iii') 映射 \(^{t}u\colon F_{b}^{\prime}\to E_{b}^{\prime}\) 是弱紧的。

\begin{enumerate}
\def\labelenumi{\arabic{enumi})}
\setcounter{enumi}{28}
\tightlist
\item
  \begin{enumerate}
  \def\labelenumii{\alph{enumii})}
  \tightlist
  \item
    证明巴拿赫空间 \(\ell^{1}(\mathbf{N})\) 的每个弱紧子集是紧的（参见 EVT, IV, p. 69, 习题 11）。特别地，从局部凸空间 E 到 \(\ell^{1}(\mathbf{N})\) 的每个弱紧线性映射是紧的。
  \end{enumerate}
\end{enumerate}

\begin{enumerate}
\def\labelenumi{\alph{enumi})}
\setcounter{enumi}{1}
\tightlist
\item
  设 E 是巴拿赫空间。由 a) 推出，从 \(c_{0}\) 到 E 的每个弱紧线性映射是紧的。
\end{enumerate}

\begin{enumerate}
\def\labelenumi{\arabic{enumi})}
\setcounter{enumi}{29}
\tightlist
\item
  \begin{enumerate}
  \def\labelenumii{\alph{enumii})}
  \tightlist
  \item
    设 E 和 F 是巴拿赫空间，\(u \in \mathcal{L}(\mathrm{E};\mathrm{F})\) 。证明下列性质等价：
  \end{enumerate}
\end{enumerate}

\begin{enumerate}
\def\labelenumi{(\roman{enumi})}
\item
  E 的每个弱紧子集在 \(u\) 下的像是（强）紧的；
\item
  E 的每个弱相对紧子集在 \(u\) 下的像在 F 中相对紧；
\item
  E 中每个弱收敛序列在 \(u\) 下的像在 F 中强收敛。
\end{enumerate}

（利用施穆良定理，EVT, IV, p. 36, 定理 2。）

如果 u 满足这些性质，就说 u 是半紧的 \(^{(3)}\) 。

\begin{enumerate}
\def\labelenumi{\alph{enumi})}
\setcounter{enumi}{1}
\item
  设 E 和 F 是巴拿赫空间。证明由半紧映射组成的集合 \(\mathcal{L}^{\mathrm{sc}}(\mathrm{E};\mathrm{F})\) 是 \(\mathcal{L}(\mathrm{E};\mathrm{F})\) 的一个闭子空间。
\item
  设 E 和 F 是巴拿赫空间，\(u \in \mathcal{L}(\mathrm{E};\mathrm{F})\) 。设 \(E_{1}\) 、\(F_{1}\) 是巴拿赫空间，\(f \in \mathcal{L}(\mathrm{E}_{1};\mathrm{E})\) 且 \(g \in \mathcal{L}(\mathrm{F};\mathrm{F}_{1})\) 。若 \(u: E \to F\) 是半紧的，则 \(g \circ u \circ f\) 亦然。
\item
  设 E 和 F 是巴拿赫空间，\(u \in \mathcal{L}(\mathrm{E};\mathrm{F})\) 。若 u 是紧的，则它是半紧的。反之，若 E 是自反的且 u 是半紧的，则 u 是紧的。
\item
  设 E 和 F 是巴拿赫空间。若 E 的对偶空间 E′ 是可数型的，则每个完全连续线性映射 E → F 都是紧的。
\item
  设 I 是无限集；\(\ell^{1}(\mathbf{I})\) 的恒同映射是半紧的，但不是弱紧的（参见 EVT, IV, p. 54, 习题 15）。g) \(\ell_{\mathrm{K}}^{1}(\mathbf{N})\) 到 \(\ell_{\mathrm{K}}^{2}(\mathbf{N})\) 的单射是半紧的但不是紧的（参见出处同上）。
\item
  设 E 是巴拿赫空间。每个连续线性映射 \(E \rightarrow \ell_{K}^{1}(\mathbf{N})\) 都是半紧的，且每个连续线性映射 \(\ell_{\mathrm{K}}^{1}(\mathbf{N}) \rightarrow \mathrm{E}\) 都是半紧的。
(3) 一些作者称它们为完全连续映射。
\item
  设 E 是巴拿赫空间。若 E 是自反的或其对偶是可数型的，则每个连续线性映射 \(\mathrm{E} \to \ell_{\mathrm{K}}^{1}(\mathbf{N})\) 都是紧的。
\end{enumerate}

\begin{enumerate}
\def\labelenumi{\arabic{enumi})}
\setcounter{enumi}{30}
\tightlist
\item
  设 E 是巴拿赫空间。
\end{enumerate}

\begin{enumerate}
\def\labelenumi{\alph{enumi})}
\tightlist
\item
  证明下列条件等价：
\end{enumerate}

\begin{enumerate}
\def\labelenumi{(\roman{enumi})}
\item
  对任何巴拿赫空间 F，从 E 到 F 的每个弱紧映射都是半紧的。
\item
  若 \((x_{n})\) 是 E 中对于拓扑 \(\sigma(\mathrm{E},\mathrm{E}^{\prime})\) 趋于 0 的序列，而 \((x_{n}^{\prime})\) 是 \(E^{\prime}\) 中对于拓扑 \(\sigma(\mathrm{E}^{\prime},\mathrm{E}^{\prime\prime})\) 趋于 0 的序列，则序列 \(\langle x_{n},x_{n}^{\prime}\rangle\) 趋于 0。
\end{enumerate}

（为证明 (i) 蕴含 (ii)，取 \(F = c_0\) 。若 (i) 不被满足，则存在从 E 到某个巴拿赫空间 F 的一个弱紧映射 \(u\) 以及 E 中弱趋于 0 的一个序列 \((x_n)\)，使得 \(\inf(\|u(x_n)\|) > 0\) ；构造一个序列 \((y_n)\)（在 \(E'\) 的单位球中），它对于 \(\sigma(E', E'')\) 收敛且满足 \(\langle u(x_n), y_n \rangle = \|u(x_n)\|\) ，并取 \(x_n' = y_n - \lim y_n\) 。）

这些条件特别在下列情形成立：\(E = \mathcal{C}(T)\)（其中 T 是紧空间），以及 \(E = L^{1}(X, \mu)\) 或 \(E = L^{\infty}(X, \mu)\)（其中 X 是局部紧空间，\(\mu\) 是 X 上的一个测度）（INT, VI, § 2, 习题 23）。b) 证明具有上述性质的巴拿赫空间不拥有无限维的自反直因子。

\begin{enumerate}
\def\labelenumi{\alph{enumi})}
\setcounter{enumi}{2}
\tightlist
\item
  设 H 是自反巴拿赫空间，B 是 H' 的赋有弱拓扑的单位球。证明 H 可以等同于 \(\mathcal{C}(\mathrm{B})\) 的一个闭的非直接子空间。
\end{enumerate}

\begin{enumerate}
\def\labelenumi{\arabic{enumi})}
\setcounter{enumi}{31}
\item
  一个拥有绍德尔基的可数型巴拿赫空间（EVT, IV, p. 70, 习题 14）具有逼近性质。
\item
  \begin{enumerate}
  \def\labelenumii{\alph{enumii})}
  \tightlist
  \item
    证明逼近空间序列的严格归纳极限是逼近空间。
  \end{enumerate}
\end{enumerate}

\begin{enumerate}
\def\labelenumi{\alph{enumi})}
\setcounter{enumi}{1}
\tightlist
\item
  证明一族逼近空间的局部凸直和是逼近空间。
\end{enumerate}

\begin{enumerate}
\def\labelenumi{\arabic{enumi})}
\setcounter{enumi}{33}
\tightlist
\item
  设 E 是巴拿赫空间。证明下列性质等价：
\end{enumerate}

\begin{enumerate}
\def\labelenumi{(\roman{enumi})}
\item
  空间 E' 是逼近空间；
\item
  对每个巴拿赫空间 F，\(\mathcal{L}^{\mathrm{f}}(\mathrm{E};\mathrm{F})\) 在 \(\mathcal{L}(\mathrm{E};\mathrm{F})\) 中的闭包是 \(\mathcal{L}^{\mathrm{c}}(\mathrm{E};\mathrm{F})\) ；
\item
  对 \(c_{0}\) 的每个闭子空间 F，\(\mathcal{L}^{\mathrm{f}}(\mathrm{E};\mathrm{F})\) 在 \(\mathcal{L}(\mathrm{E};\mathrm{F})\) 中的闭包是 \(\mathcal{L}^{\mathrm{c}}(\mathrm{E};\mathrm{F})\) 。
\end{enumerate}

\begin{enumerate}
\def\labelenumi{\arabic{enumi})}
\setcounter{enumi}{34}
\tightlist
\item
  设 E 是巴拿赫空间。我们称 E 具有度量逼近性质，如果 E 的恒同映射属于范数 \(\leqslant 1\) 的元素（\(\mathcal{L}^{\mathrm{f}}(\mathrm{E})\) 中的）组成的集合在赋有紧收敛拓扑的空间 \(\mathcal{L}(\mathrm{E})\) 中的闭包。
\end{enumerate}

证明下列性质等价：

\begin{enumerate}
\def\labelenumi{(\roman{enumi})}
\item
  空间 E 具有度量逼近性质；
\item
  \(\mathcal{L}(\mathrm{E})\) 的单位球是范数 \(\leqslant 1\) 的元素（\(\mathcal{L}^{\mathrm{f}}(\mathrm{E})\) 中的）组成的集合对于紧收敛拓扑的闭包；
\item
  对每个巴拿赫空间 F，\(\mathcal{L}(\mathrm{E};\mathrm{F})\) 的单位球是范数 \(\leqslant 1\) 的元素（\(\mathcal{L}^{\mathrm{f}}(\mathrm{E};\mathrm{F})\) 中的）组成的集合对于紧收敛拓扑的闭包；
\item
  对每个巴拿赫空间 F，\(\mathcal{L}(\mathrm{F};\mathrm{E})\) 的单位球是范数 \(\leqslant 1\) 的元素（\(\mathcal{L}^{\mathrm{f}}(\mathrm{F};\mathrm{E})\) 中的）组成的集合对于紧收敛拓扑的闭包。
\end{enumerate}

\begin{enumerate}
\def\labelenumi{\arabic{enumi})}
\setcounter{enumi}{35}
\tightlist
\item
  \begin{enumerate}
  \def\labelenumii{\alph{enumii})}
  \tightlist
  \item
    设 X 是局部紧拓扑空间。证明空间 \(\mathcal{C}_{0}(\mathrm{X})\) 具有度量逼近性质（改造 III, p. 21 命题 13 的证明）。
  \end{enumerate}
\end{enumerate}

\begin{enumerate}
\def\labelenumi{\alph{enumi})}
\setcounter{enumi}{1}
\tightlist
\item
  设 X 是赋有正测度的局部紧拓扑空间，且 \(p \in [1, +\infty]\) 。证明空间 \(L^{p}(X)\) 具有度量逼近性质（同样的方法）。
\end{enumerate}

\begin{enumerate}
\def\labelenumi{\arabic{enumi})}
\setcounter{enumi}{36}
\tightlist
\item
  设 X 是 \(C^{r}\) 类（\(1 \leqslant r \leqslant \infty\)）微分流形，仿紧且局部有限维，并设 E 是 X 上的 \(C^{r}\) 类向量丛，局部有限维。证明空间 \(\mathcal{S}^{r}(\mathrm{X};\mathrm{E})\) 具有逼近性质。
\end{enumerate}

